\documentclass[a4paper,11pt]{article}
\usepackage[T1]{fontenc}
\usepackage[utf8]{inputenc}
\usepackage[top=2.5cm,bottom=2.5cm,left=2.5cm,right=2.5cm]{geometry}
\usepackage{hyperref}
\usepackage{xcolor}
\definecolor{accent}{HTML}{1F4E78}
\hypersetup{colorlinks=true, urlcolor=accent}
\pagestyle{empty}
\setlength{\parindent}{0pt}
\setlength{\parskip}{10pt}

\begin{document}

Dear Hiring Team,

I am applying for the Product Manager role in Copenhagen. Queue-it's problem, keeping traffic fair and infrastructure safe at massive scale, is the kind of technical product problem I want to spend my time on.

At FlowSpace, a workspace marketplace, I worked as the product owner on a live platform. I led the migration off a legacy stack onto a coded Next.js and Laravel platform, mapping how each part of the system worked before deciding how to move it, while keeping the product live throughout. I integrated OfficeRnD and Nexudus for automated bookings, built support agents that answer from live platform data and escalate on low confidence, and shipped a listing pipeline that turns a provider's raw documents into a structured draft with human review. The platform grew from 2,000 to 4,500 plus users. I was also the bridge between our internal team, an external agency, and the people using the product, which meant translating customer problems into scoped work and confirming the fix actually solved them.

I bring three years of hands-on delivery across eight European companies, comfort reading architecture and APIs well enough to discuss trade-offs directly with engineers, and a habit of prototyping before committing to a build.

I hold Italian EU long-term resident status. I am not currently authorised to work in Denmark and would need the employer-side permit or sponsorship process, and I am ready to relocate and complete it. Available immediately.

Regards,\\
Prabhavit Gupta\\
+39 351 610 7807 $\vert$ prabhavitg@gmail.com $\vert$ linkedin.com/in/prabhavit

\end{document}
